\documentclass[11pt]{article}
\input{../preamble/preamble_latex.tex}
\input{../preamble/preamble_math.tex}
\title{STATS 305C: Week 2 Math Problem}
\date{Due Weds April 15, 2026 at 11:59pm}

\begin{document}
\maketitle

\noindent
Let $Y_1, \ldots, Y_n \overset{\text{i.i.d.}}{\sim} F$, where $F$ is a mixture of two
Poisson distributions: with probability $\pi$ an observation is drawn from
$\mathrm{Poi}(\mu_1)$, and with probability $1 - \pi$ it is drawn from
$\mathrm{Poi}(\mu_0)$.

\begin{enumerate}[label=(\alph*)]

    \item Write out the log-likelihood $\ell(Y;\, \mu_0, \mu_1, \pi)$. Why is direct
    maximization of this log-likelihood difficult?

    \item Introduce latent binary labels $Z_1, \ldots, Z_n$, where $Z_i = 1$ means
    observation $i$ was drawn from the $\mathrm{Poi}(\mu_1)$ component.
    \begin{enumerate}[label=(\roman*)]
        \item Specify the hierarchical model: the marginal distribution of $Z_i$
        and the conditional distribution of $Y_i \mid Z_i$.
        \item Write out the complete-data log-likelihood
        $\ell(Y, Z;\, \mu_0, \mu_1, \pi)$.
    \end{enumerate}

    \item We now use the EM algorithm to estimate the parameters in part~(a).

    \medskip
    \noindent\textit{E-step.}
    Given current estimates
    $\hat\theta^{(t)} = \bigl(\hat\mu_0^{(t)},\, \hat\mu_1^{(t)},\, \hat\pi^{(t)}\bigr)$,
    compute
    \begin{align*}
        \hat\tau_i^{(t)}
        \;=\;
        \mathbb{E}_{\hat\theta^{(t)}}[Z_i \mid Y_i]
        \;=\;
        \mathbb{P}_{\hat\theta^{(t)}}(Z_i = 1 \mid Y_i),
    \end{align*}
    then form the $Q$-function
    \begin{align*}
        Q\!\left(\theta;\, \hat\theta^{(t)}\right)
        \;=\;
        \mathbb{E}_{\hat\theta^{(t)}}\!\left[\ell(Y, Z;\, \theta) \;\middle|\; Y\right].
    \end{align*}

    \medskip
    \noindent\textit{M-step.}
    Maximize $Q$ with respect to $\theta$ and show that the closed-form update
    equations are
    \begin{align*}
        \hat\pi^{(t+1)}
        &\;=\;
        \frac{1}{n}\sum_{i=1}^n \hat\tau_i^{(t)}, \\[6pt]
        \hat\mu_0^{(t+1)}
        &\;=\;
        \frac{\displaystyle\sum_{i=1}^n \bigl(1 - \hat\tau_i^{(t)}\bigr)\,Y_i}
             {\displaystyle\sum_{i=1}^n \bigl(1 - \hat\tau_i^{(t)}\bigr)}, \\[6pt]
        \hat\mu_1^{(t+1)}
        &\;=\;
        \frac{\displaystyle\sum_{i=1}^n \hat\tau_i^{(t)}\,Y_i}
             {\displaystyle\sum_{i=1}^n \hat\tau_i^{(t)}}.
    \end{align*}

    \medskip
    \noindent\textit{Convergence.}
    What quantity should you monitor to assess convergence? Explain why monitoring
    $Q$ is not sufficient.

    \item After convergence, describe an approach to constructing an asymptotic 95\%
    confidence interval for the mixture proportion $\pi$.
    (\textit{Hint:} Observed Fisher information or Bootstrap.)

\end{enumerate}

\end{document}